\documentclass[11pt]{article}
\usepackage[T1]{fontenc}
\usepackage{lmodern,amsmath,amssymb,amsthm}
\usepackage[margin=1in]{geometry}
\usepackage[hidelinks]{hyperref}
\newtheorem{theorem}{Theorem}
\newcommand{\E}{\mathbb{E}}
\newcommand{\Prob}{\mathbb{P}}
\newcommand{\erfc}{\operatorname{erfc}}
\title{A sharp bound for Brownian search with renewal resetting\\
\large Exact clocks for a family of power-law target distances}
\author{Mathematical Physics Notes}
\date{Working paper, 26 November 2026}
\begin{document}
\maketitle
\begin{abstract}
A Brownian searcher repeatedly returns to its starting point while looking
for a target sampled once at the beginning of the search. For clocks drawn
independently from the same law at each restart, we express the duration of a
single attempt as a symmetric integral operator applied to its success
probability. An arithmetic--geometric mean inequality then bounds the mean
search time from below and characterizes equality. For a Lomax family of
target distances we construct a clock attaining the bound and compute the
minimum. Two thresholds separate divergence of the search time from
divergence of the mean interval of an optimal clock.
\end{abstract}

\section{Model and relation to previous work}
On the real line a Brownian particle starts from zero with variance $2Dt$
at time $t$, where $D>0$. A target distance $X$ has density $p$ on
$(0,\infty)$ and is sampled once. By reflection symmetry we may place the
target at $+X$. In each attempt the searcher runs a fresh Brownian motion
and draws a fresh clock $R$, independent of the motion and the target, from
a common law with $\Prob(0<R<\infty)=1$. Hitting the target ends the search.
If the clock rings first, the particle returns instantaneously to zero.
Write $T$ for the total time and average over the target, paths, and clocks.
This class includes deterministic periods and clocks with infinite mean.
Clocks updated from observations or from the number of failures are outside
the optimization domain.

Evans and Ray~\cite{evansray} derived a conjugacy relation between the
target density and the squared probability of success in one attempt.
Their variational argument gives stationarity, and their numerical
comparisons examine minimum behavior in particular clock families. They
also studied heavy-tailed clocks and conjugate target distributions.
Here a global inequality in one dimension supplies a sufficient as well as
a necessary equality condition when its bound is finite. We use it to solve
the optimization for a specified family of target densities.

\section{A universal bound and its equality condition}
Let $H_x$ denote the hitting time of $x>0$ by a fresh Brownian path, and put
\[
 q(x)=\Prob(H_x<R),\qquad a(x)=\E[\min(H_x,R)].
\]
For each $x>0$, $q(x)>0$. Hitting times have continuous laws, so ties with
an independent clock have probability zero. The expected duration of attempt
$n$, weighted by the probability of reaching it, is $(1-q(x))^{n-1}a(x)$.
Summing nonnegative terms gives $\E[T\mid X=x]=a(x)/q(x)$, also when
$a(x)=\infty$.

\begin{theorem}\label{thm:bound}
For every admissible clock and target density,
\begin{equation}\label{eq:bound}
 \E T\ \ge\ \frac{1}{D}\int_0^\infty
 \left(\int_s^\infty \sqrt{p(x)}\,dx\right)^2 ds .
\end{equation}
If $p>0$ almost everywhere and the right side is finite, equality holds
precisely for clocks satisfying $q=c\sqrt p$ almost everywhere, for some
constant $c>0$.
\end{theorem}
\begin{proof}
For a deterministic clock $t>0$, the reflection principle gives
$q_t(x)=\erfc(x/\sqrt{4Dt})$. The expected duration of an attempt is
$a_t(x)=\int_0^t(1-q_s(x))\,ds$. For $x>0$,
$\partial_s q_s=D\partial_x^2 q_s$ and $q_{0+}(x)=0$, hence
$-Da_t''(x)=q_t(x)$. Differentiation under the integral on compact
subsets of $(0,\infty)$ follows from the Gaussian decay as $s\downarrow0$.
Also $a_t(0)=0$ and $a_t'(\infty)=0$, the latter following from
$-\partial_x q_s(x)=(\pi Ds)^{-1/2}\exp[-x^2/(4Ds)]$ by dominated
convergence. Integrating the differential equation twice gives
$Da_t(x)=\int_0^\infty \min(x,y)q_t(y)\,dy$.
Averaging over $R$ and applying Tonelli's theorem gives the identity
\begin{equation}\label{eq:kernel}
 Da(x)=\int_0^\infty \min(x,y)q(y)\,dy .
\end{equation}
Consequently $D\E T$ is the integral over $x,y>0$ of the kernel
$\min(x,y)$ multiplied by $p(x)q(y)/q(x)$. The kernel is symmetric.
Averaging the integrand with its transpose and applying the
arithmetic--geometric mean inequality bounds the multiplier below by
$\sqrt{p(x)p(y)}$. The identity
$\min(x,y)=\int_0^\infty {\bf1}_{s<x}{\bf1}_{s<y}\,ds$ and Tonelli's
theorem now give~\eqref{eq:bound}.
For finite equality the nonnegative difference of the symmetrized
integrands must vanish almost everywhere. Since $\min(x,y)>0$, this means
$p(x)/q(x)^2=p(y)/q(y)^2$ for almost every pair. Fubini's theorem gives
a common constant almost everywhere. Conversely this condition gives
pointwise equality almost everywhere.
\end{proof}
For any clock, $q$ is continuous on $(0,\infty)$ and $q(0+)=1$ by
bounded convergence. Thus a positive continuous density with finite,
positive $p(0+)$ fixes the candidate uniquely as a function,
$q(x)=\sqrt{p(x)/p(0+)}$. Attainment still requires a clock with that
success probability.

\section{Exact clocks for power-law target distances}
\begin{theorem}\label{thm:lomax}
Let $L>0$, $b>1/2$, and let the target distance have density
\[
 p_b(x)=\frac{2b-1}{L}(1+x/L)^{-2b},\qquad x>0 .
\]
For $1/2<b\le3/2$ the mean search time is infinite for every admissible
clock. For $b>3/2$ its minimum over admissible clocks is
\begin{equation}\label{eq:value}
 \min \E T=\frac{L^2}{D}\frac{2b-1}{(b-1)^2(2b-3)} .
\end{equation}
To attain it, draw $K$ from a gamma distribution of shape $b$ and rate $L$
at the start of each attempt and draw $R$ conditionally exponential with
rate $DK^2$. Every clock attaining the finite minimum has infinite mean for
$3/2<b\le2$ and finite mean for $b>2$.
\end{theorem}
\begin{proof}
When $b\le1$, the inner integral in~\eqref{eq:bound} diverges. For $b>1$
it equals
$\sqrt{(2b-1)/L}\,L(1+s/L)^{1-b}/(b-1)$.
Its square is integrable in $s$ exactly when $b>3/2$. Integration gives
the right side of~\eqref{eq:value}.

Given $K=k$, the hitting-time Laplace transform yields
$\Prob(H_x<R\mid K=k)=\E e^{-Dk^2 H_x}=e^{-kx}$.
Indeed the bounded solution of $Du''=Dk^2u$ on $(-\infty,x)$ with
$u(x)=1$ and $u(-\infty)=0$ has $u(0)=e^{-kx}$. Integrating over the
gamma density $L^b k^{b-1}e^{-Lk}/\Gamma(b)$ gives
$q(x)=(1+x/L)^{-b}$. This realizes the equality condition in
Theorem~\ref{thm:bound} whenever the bound is finite.

For any clock, integration of $xq_t(x)$ gives
$\int_0^\infty xq_t(x)\,dx=Dt$, since
$\int_0^\infty z\erfc(z)\,dz=1/4$ by integration by parts. Tonelli gives
$D\E R=\int_0^\infty xq(x)\,dx$, including infinite values.
Every finite minimizer has the same normalized success function.
Substitution gives divergence when $b\le2$ and
$\E R=L^2/[D(b-1)(b-2)]$ when $b>2$.
\end{proof}

\section{A direct check at the clock threshold}
At $b=2$ the constructed clock has infinite mean but a finite optimal
mean search time. To check the value without the symmetric kernel, condition
on the exponential clock rate. One attempt then has expected duration
$(1-e^{-kx})/(Dk^2)$. Mixing against the gamma density at $b=2$ gives
\[
 a(x)=\frac{L^2}{D}\int_0^\infty
 \frac{e^{-Lk}-e^{-(L+x)k}}{k}\,dk
 =\frac{L^2}{D}\log(1+x/L).
\]
The last identity follows by integrating $e^{-uk}$ over $L<u<L+x$
first. Dividing by $q(x)$ and averaging over $p_2$ gives
\[
 \E T=\frac{3L^2}{D}\int_1^\infty \frac{\log v}{v^2}\,dv
 =\frac{3L^2}{D}.
\]
Long scheduled intervals therefore coexist with a finite mean time spent
searching: an attempt stops early whenever the target is reached.

\paragraph{Reproducibility and status.}
The accompanying \texttt{verify.py} evaluates the conditional-clock renewal
calculation independently of the lower-bound integral for four parameter
values. This working paper has not undergone peer review. Generative AI
(OpenAI Codex) assisted with the derivation, manuscript, and numerical
script. Priority for the specific bounds and thresholds remains subject
to independent literature review.

\begin{thebibliography}{1}
\bibitem{evansray} M. R. Evans and S. Ray,
Stochastic Resetting Prevails Over Sharp Restart for Broad Target
Distributions, \textit{Phys. Rev. Lett.} \textbf{134}, 247102 (2025).
\href{https://doi.org/10.1103/2wpq-gl2j}{doi:10.1103/2wpq-gl2j}.
\end{thebibliography}
\end{document}
